\documentclass{situationist-preprint}

\title{Identity and the Childhood Self}
\msid{PID-07-a}
\msfamily{PID-07}
\msdiscipline{Personal Identity and the Self}
\mstype{Principal manuscript}
\mskeywords{childhood; personal identity; psychological connectedness; obligations; promises; responsibility; Parfit}

\begin{document}
\maketitle

\begin{abstract}
Are you the same person as the child you were? We argue that you are, but that your psychological connectedness with your childhood self is weak, and that the obligations this relation grounds are correspondingly weak. Identity holds by the psychological criterion of PID-01, since overlapping chains of connections run from the child to the adult, and by the biological criterion of PID-05. Connectedness, the holding of direct connections of memory, intention, belief, desire and character, comes in degrees, and between an adult and the child of decades before it is slight. Following PID-06, on which connectedness rather than identity carries most of what matters, we argue that the force of obligations passing between a person and her childhood self tracks connectedness, with a small residue for identity. We answer the objections that this licenses breaking promises, that it excuses childhood wrongs too readily, and that it contradicts the identity claim.
\end{abstract}

\section{The problem}

Locke made memory the criterion of personal identity \citep[II.xxvii]{Locke1694}, and Reid replied with the case of a general who remembers his deeds as a young officer, who in turn remembered being flogged as a boy, though the general remembers nothing of the boy \citep[III.6]{Reid1785}. On Locke's criterion the general is the officer and the officer is the boy, but the general is not the boy, which violates the transitivity of identity. The modern psychological criterion repairs the difficulty by requiring only overlapping chains of connections, and so delivers the verdict that the general is the boy \citep{Parfit1984}. But Reid's case leaves a question that the repair does not answer. The general is the boy; what follows for what he owes to, and on account of, the boy? \citet[ch.~15]{Parfit1984} argued that the weakening of connections over a life can weaken commitments and desert, and \citet{Schechtman1996} that a person's relation to her past depends on its integration into her present self-understanding. We draw the consequence for obligations.

\section{Framework}

\begin{definition}[Connectedness; continuity]\label{d:rel}
A later person-stage is \emph{psychologically connected} with an earlier one to the degree that direct connections hold between them: memories of experiences, the carrying out of intentions, and the persistence of beliefs, desires, values and traits of character. It is \emph{psychologically continuous} with it if there is a chain of person-stages from the earlier to the later in which each is strongly connected with its successor.
\end{definition}

\begin{definition}[Childhood self]\label{d:child}
A person's \emph{childhood self} is the series of her person-stages before adolescence; we write $c$ for a representative stage of middle childhood and $a$ for a stage of mature adulthood several decades later.
\end{definition}

\begin{principle}[Weight of connectedness]\label{p:weight}
The strength of an obligation grounded in the relation between a person at one time and herself at another is an increasing function of the degree of psychological connectedness between the two, together with a small residual weight for identity.
\end{principle}

Principle~\ref{p:weight} is the deontic counterpart of the conclusion of PID-06, that connectedness carries most of what matters in survival and identity keeps a small residual weight. Its rationale is the same. An obligation grounded in a person's own past act, a promise or a wrong, binds her because the act expressed her will, and what makes a present will answerable for a past act is that the attitudes the act expressed are still in some measure hers.

\section{The argument}

\begin{lemma}\label{l:id}
In a normal life, $a$ and $c$ are stages of the same person.
\end{lemma}

\begin{proof}
In a normal life, each stage between $c$ and $a$ is strongly connected with stages a short time later: memories, intentions, beliefs and character carry over from one month to the next. So $a$ is psychologically continuous with $c$ by Definition~\ref{d:rel}, the chain does not branch, and by the criterion of PID-01 they are stages of one person. The same verdict follows from the biological criterion that PID-05 applies to the organism we are, since the organism of $c$ continues its life as the organism of $a$.
\end{proof}

\begin{lemma}\label{l:weak}
The psychological connectedness between $a$ and $c$ is weak.
\end{lemma}

\begin{proof}
Consider each kind of direct connection. Autobiographical memory of middle childhood is sparse in adulthood and largely reconstructed; memory of early childhood is almost absent. The intentions of the child, to become an astronaut or to own a horse, are rarely carried out by the adult, and when the adult's projects resemble them, it is seldom because the child's intention is executed. The child's beliefs are mostly superseded, and its values were largely taught rather than its own. Traits of temperament persist to some degree, but character in the full sense, a stable structure of reasons and commitments, is formed later. Direct connections between $a$ and $c$ are therefore few and weak.
\end{proof}

\begin{theorem}\label{t:main}
You are the same person as your childhood self, but your connectedness with it is weak, and so are the obligations grounded in the relation between you and it.
\end{theorem}

\begin{proof}
The identity claim is Lemma~\ref{l:id} and the connectedness claim is Lemma~\ref{l:weak}. By Principle~\ref{p:weight}, the strength of an obligation grounded in the relation between $a$ and $c$ is an increasing function of their connectedness, which is weak, together with the small residual weight of identity. So such obligations are weak.
\end{proof}

The theorem covers obligations running forward, from promises made and wrongs done as a child, and backward, to whatever the adult may owe to the child's hopes; both are weak in the adult, though neither is null. The companion manuscript PID-07-b examines obligations incurred in childhood in detail.

\section{Objections and replies}

\begin{objection}[Promises]
If the force of a promise fades with connectedness, people can escape their promises by changing.
\end{objection}

\begin{reply}
Principle~\ref{p:weight} concerns obligations grounded in the relation between a person and her past self, and the promises of adults are usually also grounded in the expectations of the promisee, which do not fade with the promisor's psychology. Moreover, connectedness with recent stages is strong, and one cannot weaken it at will; the change that weakens connectedness over decades is not a strategy available for escaping last week's promise.
\end{reply}

\begin{objection}[Childhood wrongs]
A man who tormented a classmate at ten is the same person at forty. If obligations to repair the wrong are weak, he is excused too easily.
\end{objection}

\begin{reply}
Weak is not null: identity keeps a residual weight, and some connections, of temperament and sometimes of attitude, may persist. Where they do, as when the adult still holds the contempt that moved the child, connectedness is greater and so is the obligation. Where they do not, our practices already agree with the verdict: juvenile records are sealed, and the reproach of a childhood fault is felt as unjust when the adult has wholly changed. PID-06-a notes that responsibility over long intervals with large psychological change is weakened in the same way.
\end{reply}

\begin{objection}[Inconsistency]
If obligations do not follow identity, the claim that the adult is the child does no work.
\end{objection}

\begin{reply}
It does two pieces of work. It explains why the adult, and no one else, is the bearer of whatever obligations survive; and it accounts for the residual weight that Principle~\ref{p:weight} assigns to identity. What it does not do is fix the strength of those obligations, which on the view of PID-06 is mostly fixed by connectedness.
\end{reply}

\section{Conclusion}

The adult is the child: overlapping chains of psychological connections and the continuing life of the organism both make it so. But the direct connections between them are few, and the obligations that pass between a person and her childhood self, forward from the child's acts and backward to the child's hopes, are weak to the same degree, though identity keeps them from vanishing.

\begin{thebibliography}{99}\small
\bibitem[Locke(1694)]{Locke1694} Locke, J. (1694). \emph{An Essay Concerning Human Understanding}, 2nd edn. Cited by book, chapter and section.
\bibitem[Parfit(1984)]{Parfit1984} Parfit, D. (1984). \emph{Reasons and Persons}. Oxford University Press.
\bibitem[Reid(1785)]{Reid1785} Reid, T. (1785). \emph{Essays on the Intellectual Powers of Man}. Cited by essay and chapter.
\bibitem[Schechtman(1996)]{Schechtman1996} Schechtman, M. (1996). \emph{The Constitution of Selves}. Cornell University Press.
\end{thebibliography}

\end{document}
